\section{Stage 6: the loop and the stopping rule}\label{sec:loop}
\sectionbrief
{stages 1--5 for one data set, and the hierarchical idea that the posterior of a parent group is the prior of its children.}
{the pipeline must decide whether the fitted laws are good enough, and if not, what to change: the degree of the factors, the structure, or the grouping of the factors.}
{closes the recursion of the original idea. Status: \emph{implemented} in \texttt{loop.py} (about $1600$ lines, committed in several steps: plan v2.9, then Phase~3). It differs from the first design (kept below as ``the original specification'') in one respect that matters: it groups \emph{factors}, not regions of the domain.}

\begin{symbols}
\begin{tabular}{@{}L{2.6cm}L{6.0cm}L{4.3cm}@{}}
$\mathcal D$ & the data set: either the gridded field, or the scattered observations $(t_i,x_i,y_i)$ & input\\
$(m,dx)$ bounds & the largest $\tau$-degree and $\xi$-degree allowed per factor at this level; start at $(3,4)$, limits $(4,8)$ & \texttt{START\_BOUNDS}, \texttt{M\_LIMIT}, \texttt{DX\_LIMIT}\\
$P_{\mathrm{conc}}$ & posterior mass needed to call $k$ or a structure ``settled'': $0.9$ & hand default \texttt{P\_CONCENTRATED}\\
held-out score & how well factors fitted on the first part of the time window explain the modes of the rest; lower is better & Sec.~\ref{sec:loop}\\
raw evidence & Gaussian evidence of the \emph{raw observations} inside the span of the factors' kernels (closed form, Sec.~\ref{sec:bayes}) & \texttt{bayes\_split.raw\_evidence}\\
$\bar\rho,\ \bar c_F,\ \bar\gamma_N$ & means over posterior samples of $\rho$, $c_F$, $\gamma_N$ & Sec.~\ref{sec:core}\\
$P(\text{split}),\ P(k)$ & posterior probability of the split, and of the number of factors & Secs.~\ref{sec:bayes}, \ref{sec:agg}\\
\end{tabular}
\end{symbols}

The loop has two halves. This section describes the first: pooling, degree growth, regrouping and recursion, which find a plausible structure. The second half, which refines that structure on the raw observations and labels it honestly, is Section~\ref{sec:refine}; most of the later fixes are there (Section~\ref{sec:fixes}).

\subsection{What the implemented loop does}
\begin{theorem}{(Outer loop, as implemented) \thy{T15} \impl}\label{thm:loop}
Given posterior samples $\{(\hat u_s,\pi_s)\}$ of the field (and, when the data are scattered, the raw observations):
\begin{enumerate}[nosep]
\item \textbf{Run and pool.} Run the core on every sample and pool the factors into the cloud $\mathcal W$ (stage 3). Samples whose factor set is rejected very strongly by the raw observations get no weight.
\item \textbf{Degree growth.} If most of the posterior mass is \textsc{misfit} or \textsc{id-ker}, raise the bounds $(m,dx)$ one step. Keep the grown level only if it (a) moves mass to \textsc{fit-compatible} and (b) predicts held-out data better (a guard against fake higher-degree fits).
\item \textbf{Regroup at fixed $k$.} If $P(k)$ is concentrated but the samples disagree on the structure, refit every sample with the consensus structure.
\item \textbf{Recursion per factor.} Analyse each factor's component field again on its own; a factor whose component still splits is replaced by its children.
\item \textbf{Decide, refine, split.} Choose the structure by the evidence of the raw data; refit the coefficients on the observations; split reducible quadratics and conjugate pairs; run the split layer (stage 5); attach intervals.
\item \textbf{Output} the factors with probabilities, intervals and a label. The domain is \emph{never} split by the loop: where each factor carries its energy is only an optional report.
\end{enumerate}
\end{theorem}
\begin{intuition}
A single run asks: given what we can see, is the law \emph{good enough}? If the data are too rich for the allowed degree, grow it, but only if the larger model also predicts data it has not seen. If the samples disagree, force them to agree and look again. If one factor hides two, split it. None of these steps moves the likelihood: they change which candidates are scored, never how they score.
\end{intuition}

\begin{explained}{the pieces the loop calls, one line each}
\textbf{\texttt{decide\_structure}}: among the candidate structures (including sub-structures), choose by the evidence of the raw observations in the kernel span, with a symbol Occam factor, a check that enough distinct wavenumbers exist, and an outlier-bin mixture. This stops the GP's own decay at high $\kappa$ from being reported as a factor.

\medskip\textbf{\texttt{screen\_samples}}: each sample's factor set is scored by the observations; samples rejected very strongly are dropped from the pool.

\medskip\textbf{\texttt{support\_path}, \texttt{refine\_factors}} (\texttt{refine\_raw}): the factor coefficients are refitted to the observations by maximising the closed-form evidence (variable projection), with terms selected by evidence minus an Occam penalty $\tfrac12p\log n$.

\medskip\textbf{\texttt{split\_quadratics}, \texttt{split\_branches}, \texttt{conjugate\_closure}}: a quadratic $\tau^2+p(\xi)\tau+q(\xi)$ is reducible over $\mathbb C$ exactly when its discriminant is a square; for a real field the annihilator is closed under complex conjugation, so complex factors are kept with their partner.

\medskip\textbf{\texttt{split\_layer}}: the Friedrichs layer of Section~\ref{sec:bayes} on every sample; \textbf{\texttt{coef\_draws}} and \textbf{\texttt{bootstrap\_intervals}}: coefficient uncertainty is propagated into the weights and, for a single noisy grid, a parametric bootstrap gives the intervals.

\medskip\textbf{\texttt{domain\_report}}: a report only, the share of each factor's energy in each block of $(t,x)$.
\end{explained}

\subsection{The original specification, and what changed}
The first design was: stop if the group is clean (at most $2$ factors, $\bar\rho<\varepsilon_1$, $\bar c_F<\varepsilon_2$, $\bar\gamma_N$ above its floor, top structure weight $\ge0.9$); otherwise split by evidence (quadtree for regions, mode space for superposed factors) and recurse, with the parent's posterior as the prior of its children. In plan v2.9 this was changed: the loop groups by \emph{factors} only (the domain tree of Section~\ref{sec:group} stays a separate tool), because grouping data points cannot separate factors and the observed failures were not grouping failures.

\begin{explained}{the loop in plain words}
\textbf{Step 1.} Run stages 1--5. \textbf{Step 2.} Look at the answer: if the posterior is concentrated on one factor set and one structure, and the held-out check passes, stop. \textbf{Step 3.} If most mass is on \textsc{misfit} or \textsc{id-ker}, raise the degree and ask the held-out guard whether the larger level is real. \textbf{Step 4.} If $k$ is settled but the samples disagree on the shape, refit all of them with the consensus shape. \textbf{Example.} Heat alone: one factor, no growth (a test checks that the loop does \emph{not} raise the degree). Transport plus diffusion: two factors by the $k$ posterior, grouped by factor. The law $\tau+0.002\,\xi^5$ needs $\xi$-degree $5$ while the start bound is $4$: at the first level no mass is \textsc{fit-compatible}, the bounds grow to $(3,5)$, the held-out guard accepts the larger level, and the loop returns one factor with coefficient $0.002$ (a test checks exactly this).
\end{explained}

\subsection{Control quantities and who sets them}
\begin{table}[!htb]
\centering\small
\begin{tabular}{@{}L{3.3cm}L{3.0cm}L{5.0cm}L{2.2cm}@{}}
\toprule
Quantity & Role in the loop & How it is set & Status\\\midrule
$P_{\mathrm{conc}}=0.9$ & settled $k$ or structure & hand default & \impl{} (hand default)\\
$(m,dx)$ start and limits & growth path & hand defaults $(3,4)$, $(4,8)$ & \impl{} (hand default)\\
$\gamma_N$ floor & cell admissibility, error bar & fixed at $0.02$ from experiments & \impl{} (hand default)\\
Held-out guard & accept a grown level & prediction of later modes & \impl{}\\
Top structure weight & structure settled & posterior $p(\mathcal M)$ & \impl{}\\
$P(\text{split})$ & do not split a group with low value & Gaussian evidence (stage 5) & \impl{}\\
Calibrated thresholds, loss-derived decisions & replace hand defaults & isotonic map and decision rule on a simulator corpus & evaluated, no-go (Section~\ref{sec:train}); the hand defaults remain\\
\bottomrule
\end{tabular}
\caption{Control quantities. ``Hand default'' means a number fixed from experiments, which the design rule says should eventually be replaced by an inferred value.}
\end{table}

\begin{caution}[title=What is still not true of the loop]
The statement ``nothing is set by hand'' is still the design intention, not the present state: $P_{\mathrm{conc}}$, the $\gamma_N$ floor, the signal-to-noise cut, the outlier cap and the degree limits are numerical defaults. A calibration layer meant to replace the decision thresholds was built and evaluated, and found to have nothing to calibrate (Section~\ref{sec:train}); a controller that would use a diagnosis to choose remedies was not built, because no available trigger finds the failures it could repair.
\end{caution}

\begin{takeaway}
The loop exists and is tested (Section~\ref{sec:limits}). It is a conservative, evidence-driven procedure whose steps change the candidates and never the likelihood. The part that is still design only is the calibrated stopping rule and the controller above it.
\end{takeaway}
